
Neural SAT solvers based on graph neural networks have demonstrated the ability to learn heuristics that reduce search time on combinatorial satisfaction problems. NeuroSAT, which uses message-passing on literal-clause bipartite graphs, achieves approximately 85\% clause satisfaction on random 3-SAT instances. However, the nature of the remaining 15\% gap has not been characterized in detail.

This work examines whether learned message-passing generates diverse violation patterns that might support subsequent refinement mechanisms. We introduce a dual-metric assessment framework that measures two independent dimensions: (1) distance diversity—the range of normalized Hamming distances from ground truth solutions, and (2) structural diversity—the range of violation entropy values across instances. Prior evaluations have reported overall satisfaction rates without analyzing the distribution of violations in near-solutions.

We test the hypothesis that baseline NeuroSAT generates heterogeneous near-solutions sufficient for basin recovery stratification. The dual-metric gate requires d/n range > 0.20 and entropy range > 2.0. Our experiments use the G4SATBench 3-SAT easy benchmark (10-40 variables, clause-to-variable ratio approximately 4.2).

Our findings show an asymmetry: the baseline architecture meets the distance diversity criterion (d/n range = 0.265) but fails the structural diversity criterion (entropy range = 1.145). This result suggests that deterministic LSTM updates with a single learned initialization point generate solutions at different quality levels while maintaining structurally similar violation patterns.
